\documentclass[11pt]{article}
\usepackage{mathblatt}
\begin{document}
\weit
\typeout{PROBE-BEGIN ebenen-e1-k2-s4-v2}
\begin{aufgabe}{\detokenize{ebenen-e1-k2-s4-v2}}
\begin{teile}
\teil Probe

\begin{ksys3}[xyz]\rpunkt{2}{0}{0}{A}\rpunkt{0}{3}{0}{B}\rpunkt{0}{0}{2}{C}\end{ksys3}
\end{teile}
\end{aufgabe}
\typeout{PROBE-END ebenen-e1-k2-s4-v2}
\typeout{PROBE-BEGIN flaecheninhalt-und-volumen-im-raum-e1-k2-s4-v1}
\begin{aufgabe}{\detokenize{flaecheninhalt-und-volumen-im-raum-e1-k2-s4-v1}}
\begin{teile}
\teil Probe

\begin{ksys3}[x1min=-1,x1max=6,x2min=-1,x2max=6,x3min=0,x3max=5] \rpunkt{1}{0}{2}{P} \rpunkt{4}{0}{2}{Q} \rpunkt{1}{1}{4}{R} \end{ksys3}
\end{teile}
\end{aufgabe}
\typeout{PROBE-END flaecheninhalt-und-volumen-im-raum-e1-k2-s4-v1}
\typeout{PROBE-BEGIN geraden-e1-k3-s4-v1}
\begin{aufgabe}{\detokenize{geraden-e1-k3-s4-v1}}
\begin{teile}
\teil Probe

\begin{ksys3}\end{ksys3}
\end{teile}
\end{aufgabe}
\typeout{PROBE-END geraden-e1-k3-s4-v1}
\typeout{PROBE-BEGIN lagebeziehungen-zone-f6-v2}
\begin{aufgabe}{\detokenize{lagebeziehungen-zone-f6-v2}}
\begin{teile}
\teil Probe

\begin{ksys3}[xyz] \rpunkt{3}{2}{4}{A} \end{ksys3}
\end{teile}
\end{aufgabe}
\typeout{PROBE-END lagebeziehungen-zone-f6-v2}
\typeout{PROBE-BEGIN punkte-und-strecken-im-koordinatensystem-e1-k1-s1-v1}
\begin{aufgabe}{\detokenize{punkte-und-strecken-im-koordinatensystem-e1-k1-s1-v1}}
\begin{teile}
\teil Probe

\begin{ksys3} \rpunkt{4}{5}{2}{C} \end{ksys3}
\end{teile}
\end{aufgabe}
\typeout{PROBE-END punkte-und-strecken-im-koordinatensystem-e1-k1-s1-v1}
\end{document}
